%----------------------------------------------------------------------------------------
% METADATA BLOCK (YAML)  — see inbox/README.md
%----------------------------------------------------------------------------------------
% uid: ""
% school: "lycee"
% level: "2eme"
% section: ["Sciences", "Informatique"]
% chapter: "Notion de Polynômes"
% title: "Vrai ou faux ?"
% difficulty: "Facile"
% tags: ["polynomes", "degre", "racines", "vrai-faux"]
% provenance: "own"
% batch: "2026-11-19-2sc-polynomes-barycentre"
%----------------------------------------------------------------------------------------

\newcommand{\R}{\mathbb{R}}
\newcommand{\vect}[1]{\overrightarrow{#1}}

% ---STATEMENT---
Répondre par vrai ou faux en justifiant (un calcul, ou un contre-exemple).
\begin{enumerate}
  \item Le polynôme $(x+1)^3-(x-1)^3$ est de degré $3$.
  \item Si $P$ et $Q$ sont deux polynômes de degré $3$, alors $P+Q$ est de degré $3$.
  \item Le réel $2$ est une racine du polynôme $x^3-3x^2+4$.
  \item Le polynôme $x^4+x^2+1$ n'a aucune racine réelle.
  \item Un polynôme de degré $2$ peut admettre trois racines distinctes.
  \item Si $a(x-1)^2+b=x^2-2x+5$ pour tout réel $x$, alors $a=1$ et $b=4$.
\end{enumerate}

% ---SOLUTION---
\begin{enumerate}[itemsep=14pt]
  \item \textbf{Faux.}
  \begin{prop}[Identités remarquables]
  \ib{$(x+1)^3=x^3+3x^2+3x+1$} \quad et \quad \ib{$(x-1)^3=x^3-3x^2+3x-1$}.
  \end{prop}
  On soustrait :
  \[(x+1)^3-(x-1)^3=x^3+3x^2+3x+1-x^3+3x^2-3x+1\]
  \[(x+1)^3-(x-1)^3=6x^2+2\]
  Les termes en $x^3$ s'éliminent : le degré est $2$.

  \item \textbf{Faux.} Contre-exemple : $P(x)=x^3$ et $Q(x)=-x^3+1$ sont de degré $3$, mais
  \[P(x)+Q(x)=1\]
  est de degré $0$.

  \item \textbf{Vrai.}
  \begin{prop}[Racine]
  $\alpha$ est une racine de $P$ si et seulement si \ib{$P(\alpha)=0$}.
  \end{prop}
  \[2^3-3\times2^2+4=8-12+4=0\]

  \item \textbf{Vrai.} Pour tout réel $x$, $x^4\geq0$ et $x^2\geq0$, donc
  \[x^4+x^2+1\geq1>0\]
  Ce polynôme ne s'annule jamais.

  \item \textbf{Faux.}
  \begin{prop}[Nombre de racines]
  Un polynôme non nul de degré $n$ a \ib{au plus $n$ racines}.
  \end{prop}
  Un polynôme de degré $2$ a donc au plus deux racines.

  \item \textbf{Vrai.} On développe le membre de gauche :
  \[a(x-1)^2+b=ax^2-2ax+a+b\]
  \begin{prop}[Identification]
  Deux polynômes sont égaux si et seulement si leurs coefficients de même degré sont égaux.
  \end{prop}
  \[
  \left\{
  \begin{aligned}
    a&=1 &&\text{(coefficient de }x^2\text{)}\\
    -2a&=-2 &&\text{(coefficient de }x\text{)}\\
    a+b&=5 &&\text{(terme constant)}
  \end{aligned}
  \right.
  \]
  Donc $a=1$ et $b=5-1=4$.
\end{enumerate}
